\subsection{A partial Dobrushin extension and a limitation of single-update contraction}
\label{resolution-dynamics:dobrushin-attempt}

This section records a new derivation on 14 November 2026. It is not
part of the previously audited revision described above. The full
Dobrushin extension remains conjectural. The positive result below
improves the pairwise sufficient constant from $1/8$ to $1/5$ by
separating three estimates that the earlier proof bounded by the same
number. It inherits the product-alphabet completion input of
Section~\ref{sec:resolution-product-alphabets}.

\begin{theorem}[An improved pairwise sufficient region]
\label{thm:dobrushin-partial-one-fifth}
Let
\[
 \mu(x)\propto\exp\left(\beta\sum_{ij}J_{ij}x_ix_j
                         +\beta\sum_i h_ix_i\right),\qquad
 K=\beta\max_i\sum_j|J_{ij}|\le\frac15.
\]
For finite parameters, $T\ge1$, and $0<\epsilon<1$, every $T$-query quantum acceptance function
has a deterministic classical simulator of mean absolute error
$\epsilon$ using $O(T^{18}\epsilon^{-28})$ individual-spin queries.
The implied constant is absolute. There is no degree bound or field
bound, and the conclusion holds uniformly after arbitrary pinnings.
This is a query bound, not a computation-time bound.
\end{theorem}

\paragraph{A Taylor representation of the update.}
Write one site's conditional mean as
$m(\eta)=\tanh(b+u(\eta))$, where
$u(\eta)=\sum_j a_j\eta_j$ and $k=\sum_j|a_j|\le1/5$.
Here $b=\beta h_i$, $a_j=\beta J_{ij}$, and lower-case $k$ in this
proof is an incident strength, not an interaction order. Put
\[
 r=\tanh b,\qquad
 \tanh(b+z)=r+\sum_{\ell\ge1}c_\ell z^\ell,\qquad
 A=\sum_{\ell\ge1}|c_\ell|k^\ell,\quad
 C=\sum_{\ell\ge1}\ell|c_\ell|k^\ell.
\]
Absolute coefficient majorization by $\tan z$ and the addition
formula for $\tanh$ give
\[
 A\le\frac{(1-r^2)\tan k}{1-|r|\tan k},\qquad
 C\le\frac{k(1-r^2)}{(\cos k-|r|\sin k)^2}.             \tag{DT1}
\]
Indeed, the majorant for $\tanh z$ is $\tan z$; expand
$(1+r\tanh z)^{-1}$ geometrically and differentiate its positive
majorant for the second bound. These series converge absolutely in
the stated range. Since $\tan k<1/3$, the first bound implies
$|r|+A\le1$: after canceling $1-|r|$, the required inequality is
$(1+2|r|)\tan k\le1$.

Choose a constant map $+1$ or $-1$ with probabilities
$(1+r-A)/2$ and $(1-r-A)/2$. With probability
$|c_\ell|k^\ell$, choose $\ell$ indices independently with
probabilities $|a_j|/k$ and return
\[
 M(\eta)=\operatorname{sign}(c_\ell)
             \prod_{v=1}^\ell\bigl(\operatorname{sign}(a_{j_v})\eta_{j_v}\bigr).
\]
For $k=0$ use the constant update directly. The pointwise mean is
exactly $m(\eta)$, so this is an exact heat-bath map. Repeated indices
cancel in pairs. Write $S(M)$ for the resulting parity support and
let $L=\ell$ for this construction, with $L=0$ on the initial constant
choices. Then $|S(M)|\le L$.

The auxiliary description uses unbounded $\ell$, but the actual
source symbol is only the signed parity map after cancellation.
There are finitely many such maps on each finite neighborhood.
Merge all descriptions of the same map, including constant maps.
Thus the source alphabet is finite, as required by the existing
product-alphabet theorem. The variable $L$ is used only to upper-bound
costs, and is not part of an oracle answer.

\paragraph{Three distinct dependency estimates.}
First, for every fixed neighbor assignment and output $y\in\{-1,1\}$,
\[
 \mathbb E[|S(M)|\mid M(\eta)=y]
 \le \mathbb E[L\mid M(\eta)=y]
 =\frac{C+y u\,\operatorname{sech}^2(b+u)}{1+y\tanh(b+u)}.
                                                               \tag{DT2}
\]
To verify the equality, use
$\mathbf1_{M=y}=(1+yM)/2$ and
$\mathbb E[LM(\eta)]=\sum_\ell\ell c_\ell u^\ell
=u\operatorname{sech}^2(b+u)$.
This signed term is the improvement over bounding the numerator
by the entire unconditioned support cost.

Set $t=yu$, $R=y\tanh b$, and
$A_k=2k/(\cos k-\sin k)^2$. The first term in (DT2) obeys
\[
 \frac{C}{1+y\tanh(b+u)}\le A_k e^{-2t}.              \tag{DT3}
\]
For completeness, after applying (DT1) the numerator factor is
\[
 \frac{1-R^2}{1+\tanh(\operatorname{arctanh}R+t)}
 =(1-R)(\cosh t+R\sinh t)e^{-t}.
\]
The derivative of $(1-R)(\cosh t+R\sinh t)$ in $R$ is
$-\cosh t+(1-2R)\sinh t<0$ for $|t|\le1/5$ and $-1\le R\le1$.
For $t\ge0$ use $3\tanh(1/5)<1$; for $t<0$ use
$|\sinh t|<\cosh t$. Its supremum is therefore its value at
$R=-1$, namely $2e^{-t}$. Also
$(\cos k-|R|\sin k)^2\ge(\cos k-\sin k)^2$, proving (DT3).

The second term in (DT2) is $t(1-y\tanh(b+u))$. If $t<0$, drop
this nonpositive term. If $t\ge0$, bound it by $2t$ and use convexity
of $A_ke^{-2t}+2t$ on $[0,k]$. Consequently
\[
 \mathbb E[|S(M)|\mid\eta,y]
 \le\max\{A_ke^{2k},\ A_k,\ A_ke^{-2k}+2k\}
 \le \rho:=\frac23 e^{2/5}<1.                       \tag{DT4}
\]
Here $1-\sin(2k)\ge1-2k$ and $k\le1/5$ bound the first two
terms by $(2/3)e^{2/5}$ and $2/3$. For the third, use
\[
 A_ke^{-2k}+2k
 \le\frac{2k}{(1-2k)(1+2k)}+2k
 \le\frac{92}{105}<\rho.
\]
The expression in the middle is increasing on $[0,1/5]$.
The strict inequality $e^{2/5}<3/2$ follows, for example, from
$\log(3/2)>2/5$ by strict convexity of $1/(1+s)$ and integration
over $0\le s\le1/2$.

Second, an unconditioned map has expected support
\[
 \mathbb E|S(M)|\le C
 \le\frac{k}{1-\sin(2k)}\le\frac13=:q.              \tag{DT5}
\]
Third, if a map is proposed repeatedly until it produces a specified
output $y$ on fixed neighbors, the expected sum of proposed support
sizes is
\[
 \frac{\mathbb E|S(M)|}{p(y\mid\eta)}
 \le\frac{2C}{1+y m(\eta)}
 \le2 A_ke^{2k}\le\frac43e^{2/5}<2.                \tag{DT6}
\]
This last quantity need not be less than one. It bounds the number
of roots supplied to the reverse exploration, not its branching rate.

\paragraph{Transfer through the existing history construction.}
Use the proper-color heat-bath schedule of
Section~\ref{resolution-dynamics:arbitrary-degree}, with the finite
map symbols just constructed. Conditional on a complete stationary
history, distinct forward maps are independent. Equation (DT4)
therefore gives, after $s_0$ whole sweeps,
\[
 \Pr[F(Z)_i\ne x_i\mid X=x]\le\rho^{s_0},\qquad
 \mathbb E[\text{forward exploration size}\mid X=x]
 \le(1-\rho)^{-1}.
\]
Both follow by the same decreasing-time path expansion as before;
complete-history conditioning is averaged only after applying the
uniform local bound. Capping a forward evaluation after $B$ nodes
adds at most $1/[B(1-\rho)]$ to its conditional error.

For posterior source access, generate the stationary history backwards
with fresh reverse maps. Each new forward-symbol request needs its
output-history spin and the neighbor-history spins used by rejection
proposals. Conditional on the complete history, (DT6) gives at most
three expected roots per request, including the output spin. For
$s$ adaptive source requests, their expected total is at most $3s$.
In the actual lazy reverse exploration, a fresh reverse map has its
unconditioned law, and (DT5) bounds its expected children by $q=1/3$.
The queue argument from (POST) consequently bounds expected physical
history work, conditional on each $x$, by
\[
 \frac{3s}{1-q}\le\frac92s.                         \tag{DT7}
\]
This again uses complete-history conditioning for rejection costs
and the exposed-information filtration for reverse exploration;
the two conditionings are not interchanged.

Choose $\delta=\epsilon^2/(288T^2)$,
$s_0=\lceil\log(1/\delta)/\log(1/\rho)\rceil$, and
$B=\lceil[\delta(1-\rho)]^{-1}\rceil$.
The conditional coupling error is at most $2\delta$, so the quantum
hybrid error is at most $\epsilon/3$. One coherent output query
uses at most $2B$ source queries. Product-alphabet simulation at error
$\epsilon/3$ therefore uses
$s=O((BT)^6\epsilon^{-15})=O(T^{18}\epsilon^{-27})$
source queries. Cap posterior work at
$\lceil27s/(2\epsilon)\rceil$; (DT7) and Markov's inequality bound
its failure probability by $\epsilon/3$ on each input. The three
errors sum to $\epsilon$, and fixing all randomness by averaging
gives the asserted deterministic tree. Pinning removes edges and
changes only fields, so the same estimates apply. This proves
Theorem~\ref{thm:dobrushin-partial-one-fifth}.

\paragraph{Why this is not the full Dobrushin theorem.}
The desired condition is
$\alpha=\max_i\sum_j\tanh(\beta|J_{ij}|)\le1-\kappa$.
There is a representation-independent obstruction to deriving it
by requiring (DT4) with some constant below one at every single-site
update and every observed output.

\begin{proposition}[A lower bound on posterior support]
\label{prop:dobrushin-posterior-obstruction}
Let $M$ be any random deterministic map, sampled independently of
the neighbor assignment, with $\Pr[M(\eta)=+1]=p(\eta)$.
Let $S(M)$ contain every coordinate on which the realized map depends.
Then
\[
 \mathbb E[|S(M)|\mid M(\eta)=+1]
 \ge\sum_j\left(1-\frac{p(\eta^{(j)})}{p(\eta)}\right)_+.
                                                               \tag{DT8}
\]
\end{proposition}
\begin{proof}
For each $j$,
$\Pr[M(\eta)=+1,M(\eta^{(j)})=-1]
\ge(p(\eta)-p(\eta^{(j)}))_+$.
On this event the map depends on $j$. Divide by $p(\eta)$ and sum.
\end{proof}

As an exact example, take a cycle with dimensionless edge strength
$a=\beta J=\tfrac12\log(3/2)$ and field
$b=\beta h=-\log8$ at every site. Its coupling-based Dobrushin
bound is $2\tanh a=2/5$. At neighbors $\eta=(+1,+1)$,
\[
 p(\eta)=\frac9{265},\qquad
 p(\eta^{(j)})=\frac1{65},\qquad
 \mathbb E[|S(M)|\mid\eta,+1]\ge\frac{128}{117}>1.   \tag{DT9}
\]
Thus even a fixed Dobrushin margin $3/5$ does not imply
single-update posterior support contraction, for any choice of
input-independent random maps. This does not contradict full AA
and does not refute the conjecture. In particular, it leaves open
block updates, multi-step estimates, weighted exploration bounds,
adaptive evaluation methods not controlled by full map support,
and direct reveal-progress arguments.

\paragraph{What can be said about a threshold in $\beta$.}
For a family with $\Lambda=\sup_n\max_i\sum_j|J_{ij}|<\infty$,
the present conclusion is the sufficient region
$\beta\Lambda\le1/5$. If $\Lambda=0$, the laws are product and
all inverse temperatures are covered. The Dobrushin conjecture would
instead cover every $\beta$ with a fixed uniform margin in
\[
 \sup_n\max_i\sum_j\tanh(\beta|J_{ij}|)<1.
\]
For maximum degree $D>1$ and equal nonzero edge magnitude $J$,
the corresponding candidate sufficient boundary is
$\operatorname{arctanh}(1/D)/J$, whereas this section proves
$\beta\le1/(5DJ)$. Neither number is established as a necessary
boundary for simulation. Forests, for example, remain simulable at
every fixed finite temperature, even beyond their Dobrushin bound.
No universal exact simulation threshold is obtained here.

Sampling the Gibbs law is a separate problem. For example,
\href{https://arxiv.org/abs/2111.04044}{Liu--Yin, \emph{Parallelize
Single-Site Dynamics up to Dobrushin Criterion}}, gives parallel
sampling results; that does not itself supply the conditional-source
oracle needed above for a particular unknown input. The remaining
research target is such an oracle or a polynomial reveal-progress
argument throughout the full Dobrushin region.

\paragraph{Verification scope.}
The algebraic estimates and the transfer argument are given explicitly
above. The accompanying script \path{../checks/dobrushin_update_checks.py}
checks finite Taylor-map examples with high-precision arithmetic and
the exact rational obstruction. These checks are supplementary, not
a proof of the dimension-uniform theorem or an independent review.
